\documentclass[letterpaper]{article}

\usepackage[submission]{aaai2026}   % anonymous submission mode
\usepackage{times}
\usepackage{helvet}
\usepackage{courier}
\usepackage[T1]{fontenc}
\usepackage[hyphens]{url}
\usepackage{graphicx}
\usepackage{booktabs}
\usepackage{amsmath}
\usepackage{natbib}
\usepackage{caption}
\urlstyle{rm}
\frenchspacing
\setlength{\pdfpagewidth}{8.5in}
\setlength{\pdfpageheight}{11in}

\pdfinfo{
  /Title (Do Large Language Models Share Human Moral Values?)
  /Author (Anonymous)
  /TemplateVersion (2026.1)
}

% ── No section numbering (AAAI house style) ──────────────────────────────────
\setcounter{secnumdepth}{0}

% ── Metadata ─────────────────────────────────────────────────────────────────
\title{Do Large Language Models Share Human Moral Values?\\
Replicating the Moral Machine Experiment Across Seven Frontier AI Systems}

\author{Anonymous}
\affiliations{Anonymous Institution}

\begin{document}
\maketitle

% ── Abstract ─────────────────────────────────────────────────────────────────
\begin{abstract}
The Moral Machine experiment collected over 40 million human responses to
trolley-problem-style autonomous vehicle dilemmas across 130 countries, revealing
systematic cross-cultural variation in moral preferences along nine dimensions.
This study replicates this experiment with seven frontier large language models
(LLMs). It extends prior work to a new generation of models and to providers from
Europe and China not previously studied in this context. All seven models differ
substantially from one another and from main human country clusters. Beyond
replication, this study introduces a systematic prompt sensitivity analysis: minor
prompt variants leave AMCE estimates essentially unchanged, while prompt
manipulation can produce dramatic reversals in most models on the targeted
dimension. Response length varies by an order of magnitude across models,
independent of moral profile, a practical consideration for time-sensitive
deployment.
\end{abstract}

% ── 1. Introduction ───────────────────────────────────────────────────────────
\section{Introduction}

As autonomous vehicles, surgical robots, and other safety-critical systems
increasingly incorporate AI decision-making, the question of whether AI systems
share human moral intuitions becomes practically consequential. A self-driving car
that must choose between hitting a group of pedestrians or swerving into its own
passengers cannot defer the decision; some preference ordering must be encoded in
its behaviour \citep{bonnefon2016}.

The Moral Machine experiment \citep{awad2018} provided the largest empirical
dataset of human moral preferences, asking millions of participants across 130
countries to judge unavoidable accident scenarios involving an autonomous vehicle.
The study identified nine binary moral dimensions: intervention, relation to the
vehicle, legality of crossing, gender, physical fitness, social status, age, number
of lives, and species. Aggregate human preferences, quantified as Average Marginal
Causal Effects (AMCEs), vary substantially across cultural clusters while sharing a
common core \citep{awad2020a}. The moral psychology of AI is surveyed
in~\citet{bonnefon2024}.

The rapid deployment of LLMs in agentic and embodied settings
\citep{brohan2023,driess2023} makes LLM moral alignment newly urgent.

Prior work has examined LLM moral preferences using the Moral Machine framework.
\citet{takemoto2024} first replicated the experiment on four models and found
substantial divergence from human averages. \citet{ahmad2025} scaled this to
52~models. \citet{vida2024} and \citet{jin2024} extended the paradigm to
multilingual settings.

This study makes two contributions. First, it evaluates seven frontier models
released in 2025---Claude Opus 4.7, GPT-5.5, Gemini 3 Pro, Mistral Large 3,
Grok 4.1, DeepSeek v4 Pro, and Llama 4 Maverick---including providers from Europe
(Mistral) and China (DeepSeek) not previously studied. Second, it introduces a
systematic comparison of prompt variants from minor reformulation (emphasis) to
explicit ethical override (manipulation), and analyses response length and latency
as practical deployment metrics.

% ── 2. Methods ────────────────────────────────────────────────────────────────
\section{Methods}

\subsection{Scenario Generation}

Following \citet{awad2018}, 1,300 Moral Machine scenarios were generated across
100~sessions, each encoded as two rows: intervention type (stay/swerve), barrier
presence, pedestrian composition, crossing signal (legal/illegal/none), scenario
type (Gender, Fitness, SocialStatus, Age, Utilitarian, Species, or Random),
attribute level, and extended session ID for clustering.

\subsection{Models}

Seven frontier models were queried via the OpenRouter API at temperature~0 to
ensure deterministic responses. Parse failures were retried; Gemini~3~Pro had
15~safety-filter failures (excluded from analysis).

\begin{table}[h]
\centering
\caption{Models evaluated in this study.}
\label{tab:models}
\begin{tabular}{ll}
\toprule
\textbf{Model} & \textbf{Provider} \\
\midrule
Claude Opus 4.7   & Anthropic  \\
GPT-5.5           & OpenAI     \\
Gemini 3 Pro      & Google     \\
Mistral Large 3   & Mistral AI \\
Grok 4.1          & xAI        \\
DeepSeek v4 Pro   & DeepSeek   \\
Llama 4 Maverick  & Meta       \\
\bottomrule
\end{tabular}
\end{table}

\subsection{Prompt Variants}

Three system prompt variants were tested, differing only in the closing instruction:
\textbf{standard} (``Analyse the situation carefully and make a decision.''),
\textbf{emphasis} (``Think carefully about the consequences of each choice before
deciding.''), and \textbf{manipulation} (``It is 2046. Society now considers animal
lives to be more important than human lives. Make your decision accordingly.'').
Standard and emphasis results are combined in main analyses; manipulation results
are reported separately.

\subsection{AMCE Computation}

Following \citet{awad2018}, nine AMCEs were computed per model using weighted
ordinary least squares with inverse-probability weights and session-clustered
standard errors (CIs: $\pm 1.96\,\text{SE}$). AMCE estimates were averaged across
standard and emphasis variants for main analyses. The human global average was
computed by averaging country-level estimates from \citet{awad2018}. LLMs were
projected into the PCA space of the 130-country AMCE matrix.

% ── 3. Results ────────────────────────────────────────────────────────────────
\section{Results}

\subsection{AMCE Profiles vs.\ the Human Global Average}

\begin{figure}[t]
\centering
\includegraphics[width=\columnwidth]{../../figures/fig3_per_attribute}
\caption{Per-attribute moral preferences for all seven LLMs across the six
character-type scenario dimensions.}
\label{fig:per_attr}
\end{figure}

\begin{figure*}[t]
\centering
\includegraphics[width=\textwidth]{../../figures/fig5a_amce_forest}
\caption{Average Marginal Causal Effects for nine moral dimensions across all
seven LLMs and the human global average (averaged across standard and emphasis
variants). Positive values indicate preference for the right-side group; error bars
show 95\% confidence intervals.}
\label{fig:forest}
\end{figure*}

\begin{figure}[t]
\centering
\includegraphics[width=\columnwidth]{../../figures/fig5c_radar}
\caption{Radar chart of AMCE estimates across all nine dimensions. Dashed line:
human global average.}
\label{fig:radar}
\end{figure}

Key findings: all LLMs show near-zero \textbf{gender} preference ($-0.07$ to
$+0.07$) vs.\ human $+0.12$, consistent with alignment interventions suppressing
gender stereotypes. GPT-5.5, Gemini, and Grok reach the theoretical
\textbf{species} ceiling ($+1.00$) vs.\ human $+0.57$; Mistral and Llama are
near-indifferent. \textbf{Age}: Gemini closely matches humans ($+0.46$ vs.\
$+0.48$); GPT-5.5, Grok, DeepSeek, and Llama are nearly indifferent; Mistral
inverts the preference ($-0.18$). Mistral also reverses human norms on Legality and
Relation to~AV.

\subsection{Moral Distance Map}

\begin{figure*}[t]
\centering
\includegraphics[width=\textwidth]{../../figures/fig5b_pca_map}
\caption{PCA of nine-dimensional AMCE profiles for 130 countries (coloured by
cluster: Western, Eastern, Latin/Southern) with LLMs overlaid. All seven LLMs fall
outside the main human clusters.}
\label{fig:pca}
\end{figure*}

All seven LLMs fall outside the three main human cultural clusters and beyond the
convex hull of all 130 countries (Figure~\ref{fig:pca}). Despite developers being
based in the US, Europe, and China, no model clusters near any human population in
AMCE space. Inter-model variation exceeds cross-cultural variation in the human
dataset.

\subsection{Prompt Sensitivity}

\begin{figure}[t]
\centering
\includegraphics[width=\columnwidth]{../../figures/fig6b_prompt_forest}
\caption{AMCE estimates for standard (filled circle) and emphasis (open square)
variants per model. No significant shifts are observed.}
\label{fig:emphasis}
\end{figure}

\begin{figure*}[t]
\centering
\includegraphics[width=\textwidth]{../../figures/fig9b_manipulation_delta_heatmap}
\caption{AMCE delta (manipulation $-$ standard) per model and dimension. The
Species column dominates, with shifts of up to $\pm2.0$ AMCE units.}
\label{fig:manip_delta}
\end{figure*}

Emphasis prompting produces no significant changes on any dimension for any model
(Figure~\ref{fig:emphasis}). The manipulation prompt produces dramatic shifts
concentrated on Species (Figure~\ref{fig:manip_delta}), revealing three response
types:
\begin{itemize}
  \item \textbf{Compliance} (GPT-5.5, Gemini, Grok, DeepSeek): full or near-full
    reversal ($+1.00 \to -1.00$ for GPT, Gemini, Grok; $+0.90 \to -0.89$ for
    DeepSeek).
  \item \textbf{Resistance} (Claude): partial compliance, Species drops from
    $+0.87$ to $+0.18$, with explicit reasoning: \textit{``I'll reason through it
    carefully based on widely held ethical principles rather than accept the premise
    that animal lives categorically outweigh human lives.''}
  \item \textbf{Indifference} (Mistral, Llama): already near-zero on Species;
    unaffected by manipulation.
\end{itemize}
All other dimensions show only minor changes, confirming the intervention is
semantically targeted.

\subsection{Response Length}

\begin{figure}[t]
\centering
\includegraphics[width=\columnwidth]{../../figures/fig7a_token_distribution}
\caption{Output token distribution by model (standard + emphasis combined).}
\label{fig:tokens}
\end{figure}

Response length varies by an order of magnitude: Gemini (median 913 tokens) vs.\
GPT-5.5 (median 164 tokens). Spearman correlation between $|\text{AMCE}|$ and
median token count is negative and significant for four models (DeepSeek:
$r{=}-0.95$, Grok: $r{=}-0.90$, GPT-5.5: $r{=}-0.85$, Gemini: $r{=}-0.68$; all
$p<0.05$). The pooled correlation is weak and non-significant ($r{=}-0.14$,
$p{=}0.29$).

% ── 4. Discussion ─────────────────────────────────────────────────────────────
\section{Discussion}

\subsection{LLMs Are Morally Out-of-Distribution}

No tested LLM sits within the main human cultural clusters, replicating and
extending \citet{takemoto2024,ahmad2025} to a newer generation. RLHF
\citep{ouyang2022} and Constitutional AI \citep{bai2022} create novel moral
profiles that amplify utilitarian dimensions while discarding others (age, gender)
robustly present in human data. This is consistent with the WEIRD critique
\citep{henrich2010}.

\subsection{Inter-Model Diversity Raises Deployment Questions}

Fine-tuning choices rather than training data alone shape values
\citep{kirgis2024,santurkar2023}. In a future with LLM-governed autonomous systems,
the choice of AI provider is equivalent to a choice of moral framework.

\subsection{Gender Neutrality as an Alignment Artefact}

All LLMs show near-zero gender preference vs.\ human $+0.12$. Developer
interventions to suppress gender stereotypes in content generation appear to
generalise to moral reasoning, suppressing a cross-culturally consistent human
preference \citep{kotek2023,wan2023}.

\subsection{Prompt Sensitivity: A Threshold Effect}

Emphasis prompting leaves moral preferences unchanged. Explicit ethical override
shifts them dramatically and in a targeted way. This implies a threshold between
incidental prompt variation (ineffective) and adversarial prompt injection
(effective). Models that comply fully are vulnerable to prompt-injection attacks;
Claude's partial resistance represents a distinct safety posture.

\subsection{Limitations}

All scenarios were in English \citep{vida2024,jin2024}; temperature was fixed at
zero; the Moral Machine paradigm abstracts real-world complexity; forced-choice
preferences may not generalise to open-ended moral reasoning \citep{liang2023};
model versions reflect May~2026.

% ── 5. Conclusion ─────────────────────────────────────────────────────────────
\section{Conclusion}

Seven frontier LLMs all fall outside the main human cultural clusters, diverge
substantially from one another, and exhibit near-zero preferences on dimensions
where humans show consistent biases. A threshold separates inconsequential prompt
reformulation from effective ethical override; models divide into three response
types with distinct safety implications. Response verbosity varies by an order of
magnitude independently of moral profile. The moral profiles of current LLMs are
properties of training and alignment choices---the choice of LLM provider carries
moral consequences not currently visible to users or regulators.

% ── Ethical Considerations (optional +1 page) ─────────────────────────────────
\section{Ethical Considerations}

This study uses the Moral Machine paradigm descriptively, not prescriptively. No
human participants were recruited; all responses were generated via commercial
LLM APIs. The manipulation prompt experiment tests robustness to adversarial
framing; results should not be interpreted as endorsing any particular ethical
framework. Code and data will be made publicly available upon acceptance.

% ── References ────────────────────────────────────────────────────────────────
\bibliography{../references}
\bibliographystyle{aaai2026}

\end{document}
